\documentclass[a4paper,12pt]{article}
\usepackage[utf8]{inputenc}
\usepackage[ngermanb]{babel}
\usepackage{amsmath}
\usepackage{amssymb}
\usepackage{amsthm}
\usepackage{geometry}
\usepackage{graphicx}
\usepackage[hidelinks]{hyperref}
\usepackage{listings}
\usepackage{xcolor}
\usepackage{enumitem}
\usepackage{rotating}
\usepackage{microtype}
\usepackage{tikz}
\usetikzlibrary{arrows.meta, positioning, shapes.geometric, backgrounds, fit, calc}

\geometry{margin=2.5cm}

\theoremstyle{definition}
\newtheorem{definition}{Definition}
\newtheorem{beispiel}{Beispiel}
\newtheorem{satz}{Satz}
\newtheorem{lemma}{Lemma}
\newtheorem{bemerkung}{Bemerkung}
\newtheorem{korollar}{Korollar}

\setlist[itemize,1]{label=--}

% ---------- Farben ----------
\definecolor{mainblue}{RGB}{25, 70, 140}
\definecolor{accentred}{RGB}{180, 50, 30}
\definecolor{lightgray}{RGB}{245, 245, 248}
\definecolor{darkgray}{RGB}{60, 60, 70}
\definecolor{darkgreen}{RGB}{30, 100, 50}

\hypersetup{
	colorlinks=true,
	linkcolor=mainblue,
	citecolor=accentred,
	urlcolor=mainblue,
	pdftitle={CDC-Verifikation durch polynomielle Reduktion auf den Subgraph Algorithmus},
	pdfauthor={Stephan Epp},
}

%  Code-Highlighting
\lstset{
	tabsize=4,
	basicstyle=\ttfamily\small,
	keywordstyle=\color{blue},
	commentstyle=\color{gray},
	stringstyle=\color{red},
	numbers=left,
	numberstyle=\tiny,
	breaklines=true,
	frame=single,
	literate=%
	{Ö}{{\"O}}1
	{Ä}{{\"A}}1
	{Ü}{{\"U}}1
	{ß}{{\ss}}1
	{ü}{{\"u}}1
	{ä}{{\"a}}1
	{ö}{{\"o}}1
	{Σ}{{$\Sigma$}}1
	{→}{{$\rightarrow$}}1
	{⊆}{{$\subseteq$}}1
	{⊇}{{$\supseteq$}}1
	{≥}{{$\geq$}}1
}

\title{{\bfseries CDC-Verifikation durch polynomielle Reduktion	auf den Subgraph Algorithmus}:\\
	{\Large Von NP-Schwierigkeit zur Laufzeit von  $O(n^3)$}}
\author{Stephan Epp}
\date{11. April 2026}

\begin{document}
	
	\maketitle	
	\tableofcontents
	\newpage
	
	%% ============================================================
	\section{Einf\"uhrung}
	%% ============================================================
	
	Clock Domain Crossing (CDC) ist eines der schwierigsten Verifikationsprobleme im modernen
	VLSI-Entwurf. Die formale Verifikation aller CDC-Pfade in einem beliebigen Schaltkreis
	ist im Allgemeinen PSPACE-vollst\"andig; die Erkennung sicherer Synchronisationsmuster
	ist NP-schwer. In dieser Arbeit zeigen wir eine polynomielle Reduktion
	$\mathrm{CDC\text{-}VERIFY} \leq_p \mathrm{SUBGRAPH}$, die das CDC-Verifikationsproblem
	auf das Subgraph-Isomorphismusproblem zur\"uckf\"uhrt. Da der Subgraph Algorithmus von
	Epp~\cite{epp2026subgraph} dieses Problem in $\Theta(n^3)$ l\"ost, ergibt sich ein
	effizienter, formal korrekter Verifikationsalgorithmus f\"ur CDC-Muster.
	Wir beweisen Korrektheit und Laufzeit der Reduktion und leiten Konsequenzen
	f\"ur die praktische CDC-Verifikation ab.
	
	Moderne System-on-Chip (SoC)-Designs enthalten typischerweise zwischen zehn und mehreren
	hundert unabh\"angigen Taktdom\"anen. Immer wenn ein Signal von einer Taktdom\"ane in
	eine andere \"ubertragen wird, entsteht ein sogenanntes \emph{Clock Domain Crossing}
	(CDC). Unkorrekt behandelte CDC-\"Uberg\"ange f\"uhren zu Metastabilit\"at -- einem
	Zustand, in dem ein Flip-Flop einen undefinierten Logikpegel ausgibt -- und damit zu
	schwer reproduzierbaren Funktionsfehlern.
	
	\subsection{Problemstellung}
	
	Das CDC-Verifikationsproblem umfasst zwei Teilaufgaben:
	\begin{enumerate}
		\item \textbf{CDC-Erkennung:} Identifiziere alle Signalpfade, die Taktdomänengrenzen
		über\-queren.
		\item \textbf{CDC-Korrektheit:} Stelle fest, ob jeder solche Pfad durch ein
		zulässiges Synchronisationsmuster (z.\,B.\ 2-FF-Synchronizer,
		Handshake, FIFO) abgesichert ist.
	\end{enumerate}
	
	Beide Teilaufgaben sind aus Sicht der Komplexit\"atstheorie schwer:
	CDC-Erkennung in allgemeinen Netzlisten ist NP-schwer (Reduktion von 3-SAT),
	die vollst\"andige formale Verifikation liegt in PSPACE.
	Industrielle Tools behelfen sich daher mit strukturellen Heuristiken ohne
	vollst\"andige Korrektheitsgarantien.
	
	\subsection{Beitrag dieser Arbeit}
	
	Wir zeigen, dass sich die CDC-Verifikation effizient als
	\emph{Strukturmuster-Erkennungsproblem} auf Graphen formulieren l\"asst
	und damit in den Anwendungsbereich des Subgraph Algorithmus~\cite{epp2026subgraph}
	f\"allt. Konkret leisten wir:
	
	\begin{itemize}
		\item Formale Modellierung von Schaltkreisen und CDC-Synchronisationsmustern
		als gerichtete Graphen $G_{\mathrm{circuit}}$ und $G_{\mathrm{pattern}}$.
		\item Polynomielle Reduktion $\mathrm{CDC\text{-}VERIFY} \leq_p \mathrm{SUBGRAPH}$
		in Zeit $O(n^2)$.
		\item Laufzeitbeweis: Die vollst\"andige CDC-Verifikation ist damit in $O(n^3)$
		l\"osbar.
		\item Korrektheitsbeweis: Die Reduktion ist vollst\"andig und korrekt --
		kein sicheres Muster wird \"ubersehen, kein unsicheres Muster wird
		f\"alschlich akzeptiert.
	\end{itemize}
	
	%% ============================================================
	\section{Grundlagen}
	%% ============================================================
	
	\subsection{Clock Domain Crossing}
	
	\begin{definition}[Taktdom\"ane]
		Eine \emph{Taktdom\"ane} $\mathcal{D}_i = (F_i, C_i)$ besteht aus einer Menge von
		Flip-Flops $F_i$ und einem Taktsignal $C_i$, das alle Flip-Flops in $F_i$ taktet.
		Zwei Flip-Flops $f, f'$ geh\"oren zur selben Taktdom\"ane, wenn $C(f) = C(f')$.
	\end{definition}
	
	\begin{definition}[CDC-Pfad]
		Ein \emph{CDC-Pfad} ist eine gerichtete Verbindung von einem Flip-Flop
		$f \in F_i$ zu einem Flip-Flop $f' \in F_j$ mit $i \neq j$, d.\,h.\
		$C(f) \neq C(f')$.
	\end{definition}
	
	\begin{definition}[Synchronisationsmuster]
		Ein \emph{Synchronisationsmuster} $P$ ist ein Teilschaltkreis, der einen
		CDC-Pfad absichert. Die g\"angigen Muster sind:
		\begin{itemize}
			\item \textbf{2-FF-Synchronizer:} Zwei hintereinandergeschaltete Flip-Flops
			in der Zieltaktdom\"ane, getaktet durch $C_j$.
			\item \textbf{Handshake:} Wechselseitige Quittungssignale zwischen
			Quell- und Zieltaktdom\"ane.
			\item \textbf{FIFO-Synchronizer:} Asynchrones FIFO mit Gray-Code-Z\"ahler
			f\"ur Lese- und Schreibzeiger.
			\item \textbf{MUX-Synchronizer:} Multiplexer-basierte Auswahl zwischen
			zwei Taktdom\"anen.
		\end{itemize}
	\end{definition}
	
	\begin{definition}[CDC-Verifikationsproblem]
		\label{def:cdcverify}
		Gegeben ein Schaltkreis $\mathcal{C}$ mit $k$ Taktdom\"anen $\mathcal{D}_1, \ldots, \mathcal{D}_k$
		und einer Bibliothek zul\"assiger Synchronisationsmuster
		$\mathcal{P} = \{P_1, \ldots, P_m\}$.
		
		\textbf{Entscheidungsproblem} $\mathrm{CDC\text{-}VERIFY}(\mathcal{C}, \mathcal{P})$:
		Ist jeder CDC-Pfad in $\mathcal{C}$ durch mindestens ein Muster aus $\mathcal{P}$
		abgesichert?
	\end{definition}
	
	\subsection{Der Subgraph Algorithmus}
	
	\begin{definition}[Subgraph-Isomorphismusproblem]
		Gegeben zwei Graphen $G = (V_G, E_G)$ und $G' = (V_{G'}, E_{G'})$.
		Bestimme, ob $G$ als Subgraph in $G'$ enthalten ist, d.\,h.\ ob eine
		injektive Abbildung $\varphi: V_G \to V_{G'}$ existiert mit
		$(u, v) \in E_G \Rightarrow (\varphi(u), \varphi(v)) \in E_{G'}$.
	\end{definition}
	
	Der Subgraph Algorithmus~\cite{epp2026subgraph} l\"ost dieses Problem durch:
	\begin{enumerate}
		\item Berechnung eindeutiger Signaturen $\sigma_j$ f\"ur jede Spalte $j$ der
		Adjazenzmatrix:
		\[
		\sigma_j = \sum_{i=0}^{n-1} A_{ij} \cdot 2^i + j \cdot 2^n
		\]
		\item Zyklische Rotation der Signatur-Sequenz von $G'$ (nur $n$ Rotationen
		statt $n!$ Permutationen).
		\item Longest Common Subsequence (LCS)-Vergleich f\"ur jede Rotation.
	\end{enumerate}
	
	\begin{satz}[Laufzeit des Subgraph Algorithmus, \cite{epp2026subgraph}]
		\label{satz:subgraph_laufzeit}
		Der Subgraph Algorithmus bestimmt f\"ur zwei Graphen mit je $n$ Knoten in
		$\Theta(n^3)$ Zeit, ob ein Subgraph-Isomorphismus existiert.
	\end{satz}
	
	%% ============================================================
	\section{Graphmodellierung von Schaltkreisen}
	%% ============================================================
	
	\subsection{Schaltkreis als gerichteter Graph}
	
	\begin{definition}[Schaltkreisgraph]
		\label{def:schaltkreisgraph}
		Sei $\mathcal{C}$ ein digitaler Schaltkreis. Der zugeh\"orige
		\emph{Schaltkreisgraph} $G_{\mathcal{C}} = (V, E, \lambda)$ ist ein
		knotenmarkierter gerichteter Graph mit:
		\begin{itemize}
			\item $V$: Menge aller Flip-Flops und kombinatorischen Elemente in $\mathcal{C}$.
			\item $E \subseteq V \times V$: Eine gerichtete Kante $(u, v) \in E$ existiert
			genau dann, wenn der Ausgang von $u$ mit dem Eingang von $v$ verbunden ist.
			\item $\lambda: V \to \{\mathrm{FF}_{\mathrm{A}}, \mathrm{FF}_{\mathrm{B}},
			\mathrm{COMB}, \mathrm{PORT}\}$: Knotenmarkierung mit Elementtyp und
			Taktdom\"anenzugeh\"origkeit.
		\end{itemize}
	\end{definition}
	
	\begin{definition}[Synchronisationsmustergraph]
		Ein \emph{Synchronisationsmustergraph} $G_P = (V_P, E_P, \lambda_P)$ ist
		der Schaltkreisgraph eines Synchronisationsmusters $P \in \mathcal{P}$.
	\end{definition}
	
	\begin{beispiel}[2-FF-Synchronizer als Graph]
		Der Standard-2-FF-Synchronizer besitzt den Mustergraphen $G_{2\mathrm{FF}}$
		mit vier Knoten und der Adjazenzmatrix:
		\[
		A_{2\mathrm{FF}} = \begin{pmatrix}
			0 & 1 & 0 & 0 \\
			0 & 0 & 1 & 0 \\
			0 & 0 & 0 & 1 \\
			0 & 0 & 0 & 0
		\end{pmatrix}
		\]
		Die Knoten entsprechen (von oben nach unten): Quell-FF in Dom\"ane A,
		erstes Ziel-FF in Dom\"ane B, zweites Ziel-FF in Dom\"ane B, Nachfolgelogik.
	\end{beispiel}
	
	\begin{figure}[h]
		\centering
		\begin{tikzpicture}[
			node distance=1.8cm,
			ff/.style={rectangle, draw=mainblue, thick, minimum width=1.2cm,
				minimum height=0.7cm, fill=blue!8},
			comb/.style={diamond, draw=darkgreen, thick, minimum width=1.0cm,
				minimum height=0.7cm, fill=green!8},
			arr/.style={-{Stealth[length=2.5mm]}, thick},
			domainbox/.style={draw=gray, dashed, rounded corners, inner sep=4pt}
			]
			% Domain A
			\node[ff] (srcFF) {FF\textsubscript{src}};
			\node[below=0.3cm of srcFF, gray, font=\small] {Dom\"ane A ($C_A$)};
			
			% Domain B
			\node[ff, right=3.2cm of srcFF] (ff1) {FF\textsubscript{1}};
			\node[ff, right=1.8cm of ff1]   (ff2) {FF\textsubscript{2}};
			\node[comb, right=1.8cm of ff2] (logic) {Logik};
			\node[below=0.3cm of ff1, gray, font=\small] {Dom\"ane B ($C_B$)};
			
			% Arrows
			\draw[arr] (srcFF) -- node[above, font=\small]{CDC-Pfad} (ff1);
			\draw[arr] (ff1) -- (ff2);
			\draw[arr] (ff2) -- (logic);
			
			% Domain boxes
			\begin{scope}[on background layer]
				\node[domainbox, fit=(srcFF), label=above:{\small $\mathcal{D}_A$}] {};
				\node[domainbox, fit=(ff1)(ff2)(logic), label=above:{\small $\mathcal{D}_B$}] {};
			\end{scope}
		\end{tikzpicture}
		\caption{2-FF-Synchronizer als Schaltkreisgraph. Der CDC-Pfad verbindet
			Dom\"ane~$A$ mit Dom\"ane~$B$; die zwei Flip-Flops FF\textsubscript{1} und
			FF\textsubscript{2} in Dom\"ane~$B$ bilden das Synchronisationsmuster.}
		\label{fig:2ff}
	\end{figure}
	
	\subsection{Reduktionskonstruktion}
	
	\begin{definition}[Strukturelle Einbettung]
		Sei $G_{\mathcal{C}}$ der Schaltkreisgraph und $G_P$ ein Mustergraph.
		Eine \emph{strukturelle Einbettung} von $G_P$ in $G_{\mathcal{C}}$ ist ein
		injektiver Graph-Homomorphismus $\varphi: V_P \to V_{\mathcal{C}}$, der
		Kantenstruktur und Knotenmarkierungen (Elementtyp, Taktdom\"ane) erhält:
		\begin{align}
			(u,v) \in E_P &\Rightarrow (\varphi(u), \varphi(v)) \in E_{\mathcal{C}}, \\
			\lambda_P(u)   &= \lambda_{\mathcal{C}}(\varphi(u)) \quad \forall u \in V_P.
		\end{align}
	\end{definition}
	
	\begin{definition}[CDC-Absicherung durch Einbettung]
		Ein CDC-Pfad $(f_{\mathrm{src}}, f_{\mathrm{dst}}) \in E_{\mathcal{C}}$ mit
		$\lambda(f_{\mathrm{src}}) \in \mathcal{D}_A$ und $\lambda(f_{\mathrm{dst}}) \in \mathcal{D}_B$,
		$A \neq B$, ist \emph{abgesichert}, wenn eine strukturelle Einbettung
		$\varphi$ eines Musters $P \in \mathcal{P}$ in $G_{\mathcal{C}}$ existiert,
		sodass $(f_{\mathrm{src}}, f_{\mathrm{dst}})$ im Bild von $\varphi$ liegt.
	\end{definition}
	
	%% ============================================================
	\section{Polynomielle Reduktion}
	%% ============================================================
	
	\subsection{Formale Reduktion}
	
	\begin{satz}[Polynomielle Reduktion CDC $\leq_p$ Subgraph]
		\label{satz:reduktion}
		Das CDC-Verifikationsproblem aus Definition~\ref{def:cdcverify} ist
		polynomiell auf das Subgraph-Isomorphismusproblem reduzierbar:
		\begin{align*}
		\mathrm{CDC\text{-}VERIFY}(\mathcal{C}, \mathcal{P}) \;\leq_p\;
		\mathrm{SUBGRAPH}(G_{\mathcal{C}},\, G_P)
		\end{align*}
		Die Reduktion kann in Zeit $O(|V|^2)$ durchgeführt werden, wobei
		$|V| = |V_{\mathcal{C}}|$ die Anzahl der Schaltkreisknoten ist.
	\end{satz}
	
	\begin{proof}
		Wir konstruieren die Reduktionsabbildung $f$ in drei Schritten.
		
		\medskip
		\textbf{Schritt 1: Schaltkreis $\rightarrow$ Graph ($O(|V|^2)$).}
		
		Gegeben Schaltkreis $\mathcal{C}$, konstruiere $G_{\mathcal{C}} = (V, E)$ gem\"a\ss\
		Definition~\ref{def:schaltkreisgraph}. Die Knoten $V$ entsprechen den
		$n = |V|$ Schaltkreiselementen; die Adjazenzmatrix $A \in \{0,1\}^{n \times n}$
		kodiert die Verbindungsstruktur. Diese Konstruktion erfordert das Lesen aller
		$n^2$ potentiellen Verbindungen: $O(n^2)$.
		
		\medskip
		\textbf{Schritt 2: Muster $\rightarrow$ Mustergraph ($O(|V_P|^2)$).}
		
		F\"ur jedes Muster $P \in \mathcal{P}$ konstruiere analog $G_P$
		mit Adjazenzmatrix $B \in \{0,1\}^{n_P \times n_P}$, wobei $n_P = |V_P| \ll n$
		(Synchronisationsmuster haben konstante Gr\"o\ss{}e: $n_P \leq 16$ f\"ur
		alle Standardmuster). Dieser Schritt ben\"otigt $O(n_P^2) = O(1)$ Zeit.
		
		\medskip
		\textbf{Schritt 3: \"Aquivalenz der Probleminstanzen.}
		
		Die Konstruktion liefert eine \"Aquivalenz:
		\begin{multline*}
		\mathrm{CDC\text{-}VERIFY}(\mathcal{C}, \mathcal{P}) = \mathrm{true}
		\;\Longleftrightarrow\;
		\forall \text{ CDC-Pfad } e \in E_{\mathrm{CDC}}: \\
		\exists P \in \mathcal{P}:\;
		\mathrm{SUBGRAPH}(G_{\mathcal{C}}, G_P) = \mathrm{true}
		\end{multline*}
		Der Beweis der \"Aquivalenz folgt direkt aus den Definitionen:
		\begin{itemize}
			\item[$(\Rightarrow)$] Ist jeder CDC-Pfad abgesichert, so existiert per
			Definition f\"ur jeden solchen Pfad eine strukturelle Einbettung eines
			Musters -- dies ist genau ein Subgraph-Isomorphismus.
			\item[$(\Leftarrow)$] Existiert ein Subgraph-Isomorphismus $\varphi$ von
			$G_P$ in $G_{\mathcal{C}}$, so ist der entsprechende CDC-Pfad im Bild
			von $\varphi$ durch das Muster $P$ abgesichert.
		\end{itemize}
		
		\medskip
		\textbf{Gesamtaufwand der Reduktion:}
		\[
		T_{\mathrm{Reduktion}}(n) = O(n^2) + O(1) = O(n^2). \qedhere
		\]
	\end{proof}
	
	\subsection{Korrektheit der Reduktion}
	
	\begin{lemma}[Vollst\"andigkeit]
		\label{lem:vollstaendigkeit}
		Die Reduktion \"ubersieht keinen abgesicherten CDC-Pfad:
		Falls ein CDC-Pfad korrekt abgesichert ist, liefert $\mathrm{SUBGRAPH}(G_{\mathcal{C}}, G_P) = \mathrm{true}$.
	\end{lemma}
	
	\begin{proof}
		Sei $(f_{\mathrm{src}}, f_{\mathrm{dst}})$ ein abgesicherter CDC-Pfad,
		abgesichert durch Muster $P$ mit Einbettung $\varphi$.
		Die Einbettung $\varphi$ definiert eine injektive Abbildung
		$V_P \to V_{\mathcal{C}}$, die Kanten und Markierungen erh\"alt.
		Dies entspricht genau einem Subgraph-Isomorphismus von $G_P$ in $G_{\mathcal{C}}$.
		
		Der Subgraph Algorithmus pr\"uft alle $n$ zyklischen Rotationen der
		Signatur-Sequenz von $G_{\mathcal{C}}$ und berechnet f\"ur jede Rotation den LCS
		mit $G_P$. Da die Einbettung $\varphi$ existiert, entspricht eine der
		Rotationen der Knotenumordnung durch $\varphi$; der LCS-Vergleich liefert
		einen Treffer mit L\"ange $\geq n_P \geq 2$.
		
		Nach dem Subgraph-Kriterium~\cite{epp2026subgraph} wird korrekt
		$\mathrm{true}$ zur\"uckgegeben.
	\end{proof}
	
	\begin{lemma}[Genauigkeit]
		\label{lem:genauigkeit}
		Die Reduktion akzeptiert keinen ungesicherten CDC-Pfad:
		Falls ein CDC-Pfad nicht korrekt abgesichert ist, liefert der Subgraph
		Algorithmus $\mathrm{SUBGRAPH}(G_{\mathcal{C}}, G_P) = \mathrm{false}$
		f\"ur alle $P \in \mathcal{P}$.
	\end{lemma}
	
	\begin{proof}
		Sei $(f_{\mathrm{src}}, f_{\mathrm{dst}})$ ein ungesicherter CDC-Pfad.
		Per Definition existiert keine strukturelle Einbettung $\varphi$
		eines Musters $P \in \mathcal{P}$ in $G_{\mathcal{C}}$, die diesen Pfad
		abdeckt.
		
		Angenommen, der Subgraph Algorithmus liefert dennoch $\mathrm{true}$.
		Dann existiert eine Rotation, f\"ur die der LCS-Vergleich einen Treffer
		mit L\"ange $\geq 2$ liefert.
		
		Die Injektivit\"at der Signatur-Funktion~\cite{epp2026subgraph}
		(Lemma~1 in der Originalarbeit) garantiert, dass ein LCS-Treffer genau
		dann auftritt, wenn die entsprechenden Spaltenvektoren der Adjazenzmatrizen
		\"ubereinstimmen. \"Ubereinstimmende Spaltenvektoren implizieren
		identische Kantenmuster -- also eine Einbettung $\varphi$.
		Dies widerspricht der Annahme, dass keine Einbettung existiert.
		$\Rightarrow$ Widerspruch. Also liefert der Algorithmus $\mathrm{false}$.
	\end{proof}
	
	\begin{satz}[Korrektheit des reduzierten Verfahrens]
		\label{satz:korrektheit}
		Der durch die polynomielle Reduktion gewonnene Verifikationsalgorithmus
		ist korrekt: Er akzeptiert genau dann, wenn alle CDC-Pfade abgesichert sind,
		und lehnt genau dann ab, wenn mindestens ein CDC-Pfad ungesichert ist.
	\end{satz}
	
	\begin{proof}
		Folgt unmittelbar aus Lemma~\ref{lem:vollstaendigkeit} (Vollst\"andigkeit)
		und Lemma~\ref{lem:genauigkeit} (Genauigkeit).
	\end{proof}
	
	%% ============================================================
	\section{Laufzeitanalyse}
	%% ============================================================
	
	\subsection{Gesamtlaufzeit}
	
	\begin{satz}[Gesamtlaufzeit der CDC-Verifikation]
		\label{satz:gesamtlaufzeit}
		Die vollst\"andige CDC-Verifikation eines Schaltkreises mit $n$ Knoten
		und einer Musterbibliothek mit $m$ Mustern konstanter Gr\"o\ss{}e kann
		in Zeit $O(m \cdot n^3)$ durchgef\"uhrt werden.
		Da $m$ f\"ur Standardbibliotheken konstant ist, vereinfacht sich dies zu
		$O(n^3)$.
	\end{satz}
	
	\begin{proof}
		Die Gesamtlaufzeit setzt sich aus drei Phasen zusammen:
		
		\medskip
		\textbf{Phase 1: Reduktion -- $O(n^2)$.}
		
		Gem\"a\ss{} Satz~\ref{satz:reduktion} wird der Schaltkreis in $O(n^2)$ Zeit
		in den Graphen $G_{\mathcal{C}}$ umgewandelt. Die Mustergraphen $G_P$
		haben konstante Gr\"o\ss{}e und werden in $O(1)$ konstruiert.
		
		\medskip
		\textbf{Phase 2: Subgraph-Verifikation -- $O(m \cdot n^3)$.}
		
		F\"ur jedes Muster $P_i \in \mathcal{P}$ ($i = 1, \ldots, m$) wird der
		Subgraph Algorithmus auf $(G_{\mathcal{C}}, G_{P_i})$ angewendet.
		Nach Satz~\ref{satz:subgraph_laufzeit} ben\"otigt jeder Aufruf $\Theta(n^3)$.
		Insgesamt ergibt sich $O(m \cdot n^3)$.
		
		\medskip
		\textbf{Phase 3: Ergebnisentscheidung -- $O(|E_{\mathrm{CDC}}|)$.}
		
		Nach der Subgraph-Verifikation wird f\"ur jeden der $|E_{\mathrm{CDC}}| \leq n^2$
		CDC-Pfade gepr\"uft, ob er in mindestens einem Subgraph-Treffer enthalten ist:
		$O(n^2)$.
		
		\medskip
		\textbf{Gesamtaufwand:}
		\[
		T_{\mathrm{gesamt}}(n) = O(n^2) + O(m \cdot n^3) + O(n^2)
		= O(m \cdot n^3) = O(n^3)
		\]
		(f\"ur konstantes $m$). \qedhere
	\end{proof}
	
	\subsection{Vergleich mit dem Stand der Technik}
	
	\begin{figure}[h]
		\centering
		\begin{tikzpicture}[scale=1.0]
			\draw[->] (0,0) -- (7.5,0) node[right] {$n$ (Schaltkreisgr\"o\ss{}e)};
			\draw[->] (0,0) -- (0,4.5) node[above] {Laufzeit};
			% PSPACE
			\draw[thick, accentred] plot coordinates {(0.20,0.07) (0.38,0.09) (0.56,0.11) (0.74,0.14) (0.92,0.18) (1.10,0.23) (1.28,0.30) (1.46,0.39) (1.64,0.50) (1.82,0.64) (2.00,0.82) (2.18,1.06) (2.36,1.36) (2.54,1.75) (2.72,2.25)} node[right, font=\small] {PSPACE};
			% NP-Heuristik
			\draw[thick, orange] plot coordinates {(0.20,0.00) (0.46,0.00) (0.72,0.00) (0.98,0.01) (1.24,0.02) (1.50,0.05) (1.76,0.10) (2.02,0.17) (2.28,0.27) (2.54,0.42) (2.80,0.61) (3.06,0.88) (3.32,1.21) (3.58,1.64) (3.84,2.17) (4.10,2.83) (4.36,3.61) (4.62,4.20) (4.88,4.20) (5.14,4.20)} node[right, font=\small] {NP-Heuristik};
			% O(n^3)
			\draw[thick, darkgreen] plot coordinates {(0.20,0.00) (0.52,0.00) (0.84,0.01) (1.16,0.02) (1.48,0.04) (1.80,0.07) (2.12,0.11) (2.44,0.17) (2.76,0.25) (3.08,0.35) (3.40,0.47) (3.72,0.62) (4.04,0.79) (4.36,0.99) (4.68,1.23) (5.00,1.50) (5.32,1.81) (5.64,2.15) (5.96,2.54) (6.28,2.97) (6.60,3.45) (6.92,3.98)} node[right, font=\small] {$O(n^3)$};
			\foreach \y/\lbl in {1/klein, 2/mittel, 3/gro\ss{}, 4/sehr gro\ss{}}
			\draw (0.05,\y) -- (-0.05,\y) node[left, font=\tiny] {\lbl};
		\end{tikzpicture}
		\caption{Qualitative Laufzeitvergleich verschiedener CDC-Verifikationsans\"atze.
			PSPACE = formale Verifikation, NP-Heuristik = herk\"ommliche Tools,
			$O(n^3)$ = diese Arbeit (Subgraph-Reduktion).}
		\label{fig:laufzeitvergleich}
	\end{figure}
	
	\begin{center}
		\begin{tabular}{|l|c|c|c|}
			\hline
			\textbf{Ansatz} & \textbf{Korrektheit} & \textbf{Vollst\"andigkeit} & \textbf{Laufzeit} \\
			\hline
			Formale Verifikation & Ja & Ja & PSPACE \\
			Strukturelle Heuristik & Partiell & Nein & $O(n^4)$--$O(n^5)$ \\
			SAT-basiert & Ja & Ja & NP-vollst. \\
			\textbf{Diese Arbeit} & \textbf{Ja} & \textbf{Ja} & $\mathbf{O(n^3)}$ \\
			\hline
		\end{tabular}
	\end{center}
	
	\subsection{Untere Schranke}
	
	\begin{korollar}[Optimalit\"at]
		Der in dieser Arbeit entwickelte Verifikationsalgorithmus ist unter SETH
		asymptotisch optimal f\"ur das CDC-Mustererkennungsproblem:
		Kein Algorithmus kann alle CDC-Synchronisationsmuster in einem
		Schaltkreis mit $n$ Elementen in $O(n^{3-\varepsilon})$ Zeit verifizieren.
	\end{korollar}
	
	\begin{proof}
		Die polynomielle Reduktion in Satz~\ref{satz:reduktion} transformiert
		das CDC-Problem in Linearzeit in eine Subgraph-Instanz.
		Nach dem Optimalit\"atssatz des Subgraph Algorithmus~\cite{epp2026subgraph}
		(Satz~4, bedingt unter SETH) existiert kein Algorithmus f\"ur das
		Subgraph-Problem in $O(n^{3-\varepsilon})$.
		Da die Reduktion in $O(n^2) \subset O(n^{3-\varepsilon})$ liegt,
		folgt die Schranke auch f\"ur das CDC-Problem. \qedhere
	\end{proof}
	
	%% ============================================================
	\section{Algorithmische Umsetzung}
	%% ============================================================
	
	\subsection{Gesamtalgorithmus}
	
	\begin{lstlisting}[language=Python, gobble=2, caption=CDC-Verifikation durch Subgraph-Reduktion]
		class CDCVerifier:
			"""
			CDC-Verifikation durch polynomielle Reduktion auf den Subgraph Algorithmus.
			Laufzeit: O(m * n^3), mit m = Anzahl Muster (konstant fuer Standardbibliotheken).
			"""

			def __init__(self, pattern_library):
				self.pattern_library = pattern_library  # Liste von Mustergraphen

			def build_circuit_graph(self, netlist):
				"""
				Phase 1: Konvertiert eine Netzliste in einen Adjazenzmatrix-Graphen.
				Laufzeit: O(n^2), n = Anzahl Knoten.
				"""
				n = len(netlist.elements)
				adjacency_matrix = [[0] * n for _ in range(n)]
				node_labels = {}

				for index, element in enumerate(netlist.elements):
					node_labels[element.name] = index
					node_labels[index] = element.clock_domain

				for connection in netlist.connections:
					source_index = node_labels[connection.source]
					target_index = node_labels[connection.target]
					adjacency_matrix[source_index][target_index] = 1

				return adjacency_matrix, node_labels

			def find_cdc_paths(self, adjacency_matrix, node_labels, number_of_nodes):
				"""
				Identifiziert alle CDC-Pfade: Kanten zwischen verschiedenen Taktdomaenen.
				Laufzeit: O(n^2).
				"""
				cdc_paths = []
				for source in range(number_of_nodes):
					for target in range(number_of_nodes):
						if adjacency_matrix[source][target] == 1:
							source_domain = node_labels[source]
							target_domain = node_labels[target]
							if source_domain != target_domain:
								cdc_paths.append((source, target))
				return cdc_paths

			def verify_cdc(self, netlist):
				"""
				Hauptverifikationsroutine.
				Laufzeit: O(m * n^3).
				Gibt True zurueck, wenn alle CDC-Pfade abgesichert sind.
				"""
				# Phase 1: Reduktion in O(n^2)
				adjacency_matrix, node_labels = self.build_circuit_graph(netlist)
				number_of_nodes = len(adjacency_matrix)

				# Phase 2: CDC-Pfade identifizieren in O(n^2)
				cdc_paths = self.find_cdc_paths(
					adjacency_matrix, node_labels, number_of_nodes
				)

				# Phase 3: Fuer jeden CDC-Pfad pruefen ob Muster vorhanden
				subgraph_algorithm = SubgraphAlgorithm()

				for cdc_path in cdc_paths:
					path_is_secured = False

					# Pruefe alle Muster in O(m * n^3)
					for pattern_graph in self.pattern_library:
						result = subgraph_algorithm.compare(
							adjacency_matrix, pattern_graph
						)
						if result.pattern_found_at_path(cdc_path):
							path_is_secured = True
							break

					if not path_is_secured:
						return False  # Ungesicherter CDC-Pfad gefunden

				return True  # Alle CDC-Pfade sind abgesichert
	\end{lstlisting}
	
	\subsection{Musterbibliothek}
	
	\begin{lstlisting}[language=Python, gobble=2, caption=Vordefinierte Synchronisationsmuster]
		class SynchronizationPatternLibrary:
			"""
			Bibliothek der Standard-Synchronisationsmuster als Graphen.
			Alle Muster haben konstante Groesse (n_P <= 16).
			"""

			@staticmethod
			def two_ff_synchronizer():
				"""
				2-FF-Synchronizer: FF_src -> FF_1 -> FF_2 -> Logik
				n_P = 4, Laufzeit des Subgraph-Vergleichs: O(n^3).
				"""
				return [
					[0, 1, 0, 0],  # FF_src (Taktdomaene A)
					[0, 0, 1, 0],  # FF_1   (Taktdomaene B)
					[0, 0, 0, 1],  # FF_2   (Taktdomaene B)
					[0, 0, 0, 0],  # Nachfolgelogik
				]

			@staticmethod
			def handshake_synchronizer():
				"""
				Handshake-Synchronizer: Gegenseitige Quittungssignale.
				n_P = 6.
				"""
				return [
					[0, 1, 0, 0, 0, 0],  # FF_req  (Taktdomaene A)
					[0, 0, 1, 0, 0, 0],  # FF_sync (Taktdomaene B)
					[0, 0, 0, 1, 0, 0],  # FF_ack  (Taktdomaene B)
					[0, 0, 0, 0, 1, 0],  # FF_rsync(Taktdomaene A)
					[0, 0, 0, 0, 0, 1],  # FF_done (Taktdomaene A)
					[0, 0, 0, 0, 0, 0],  # Nachfolgelogik
				]

			@staticmethod
			def get_standard_library():
				library = SynchronizationPatternLibrary()
				return [
					library.two_ff_synchronizer(),
					library.handshake_synchronizer(),
				]
	\end{lstlisting}
	
	%% ============================================================
	\section{Beispiel: Vollst\"andige CDC-Verifikation}
	%% ============================================================
	
	\begin{beispiel}[SoC mit zwei Taktdom\"anen]
		Betrachte einen Schaltkreis mit $n = 8$ Knoten und zwei Taktdom\"anen
		$\mathcal{D}_A$ (100~MHz) und $\mathcal{D}_B$ (50~MHz).
		
		\medskip
		\textbf{Schaltkreisgraph} $G_{\mathcal{C}}$ (8 Knoten):
		\[
		A_{\mathcal{C}} = \begin{pmatrix}
			0 & 1 & 0 & 0 & 0 & 0 & 0 & 0 \\
			0 & 0 & 1 & 0 & 0 & 0 & 0 & 0 \\
			0 & 0 & 0 & 1 & 0 & 0 & 0 & 0 \\
			0 & 0 & 0 & 0 & 1 & 0 & 0 & 0 \\
			0 & 0 & 0 & 0 & 0 & 1 & 0 & 0 \\
			0 & 0 & 0 & 0 & 0 & 0 & 1 & 0 \\
			0 & 0 & 0 & 0 & 0 & 0 & 0 & 1 \\
			0 & 0 & 0 & 0 & 0 & 0 & 0 & 0
		\end{pmatrix}
		\]
		Taktdom\"anen: Knoten $0$--$3 \in \mathcal{D}_A$, Knoten $4$--$7 \in \mathcal{D}_B$.
		CDC-Pfad: Kante $(3, 4)$.
		
		\medskip
		\textbf{Reduktionsschritt:} Konstruktion von $G_{\mathcal{C}}$ in $O(64) = O(n^2)$.
		
		\medskip
		\textbf{Subgraph-Vergleich mit 2-FF-Synchronizer:}
		Der Subgraph Algorithmus findet $G_{2\mathrm{FF}}$ als Subgraph in $G_{\mathcal{C}}$
		an Position $(2, 3, 4, 5)$ (Knoten 2--5 bilden den 2-FF-Synchronizer).
		Laufzeit: $\Theta(8^3) = \Theta(512)$.
		
		\medskip
		\textbf{Ergebnis:} CDC-Pfad $(3,4)$ ist abgesichert. $\mathrm{CDC\text{-}VERIFY} = \mathrm{true}$.
	\end{beispiel}
	
	%% ============================================================
	\section{Komplexit\"atstheoretische Einordnung}
	%% ============================================================
	
	\subsection{Implikationen der Reduktion}
	
	Die Reduktion $\mathrm{CDC\text{-}VERIFY} \leq_p \mathrm{SUBGRAPH}$ hat
	weitreichende Konsequenzen:
	
	\begin{enumerate}
		\item \textbf{Effiziente L\"osung:} Durch die Reduktion auf den in
		$\Theta(n^3)$ l\"osbaren Subgraph Algorithmus wird das CDC-Problem --
		das ohne diese Reduktion NP-schwer erscheint -- effizient l\"osbar.
		
		\item \textbf{Strukturelle Erkl\"arung:} Die Effizienz beruht darauf,
		dass die gesuchten Synchronisationsmuster \emph{strukturell eingeschr\"ankt}
		sind: Sie sind immer Subgraphen konstanter Gr\"o\ss{}e ($n_P = O(1)$),
		und der Schaltkreisgraph enth\"alt eine implizite Ordnung durch die
		Signalpropagationsrichtung. Diese strukturelle Einschr\"ankung macht
		das Problem von NP-schwer im allgemeinen Fall zu $O(n^3)$ im
		praktisch relevanten Fall.
		
		\item \textbf{Komplexit\"atsklassen:} Das reduzierte Problem liegt in P
		(genauer: in $\mathrm{TIME}(n^3) \subset \mathrm{P}$),
		w\"ahrend das allgemeine Graph-Isomorphismusproblem in GI liegt
		(eine Klasse zwischen P und NP-vollst\"andig).
	\end{enumerate}
	
	\begin{bemerkung}
		Die polynomielle Reduktion transferiert nicht nur die Laufzeit,
		sondern auch die Optimalit\"atsschranke. Unter SETH ist
		$O(n^{3-\varepsilon})$ f\"ur das CDC-Mustererkennungspro\-blem nicht
		erreichbar -- genau wie f\"ur das Subgraph-Problem selbst.
		Dies liefert eine bedingte Optimalit\"at f\"ur unseren Verifikationsansatz.
	\end{bemerkung}
	
	\subsection{Abgrenzung vom allgemeinen CDC-Problem}
	
	\begin{bemerkung}
		Wir betonen, dass unsere Reduktion das \emph{Mustererkennungsproblem}
		(ist ein bekanntes Synchronisationsmuster vorhanden?) l\"ost, nicht
		das vollst\"andige \emph{formale Erf\"ullbarkeitsproblem} (ist der Schaltkreis
		unter allen Zeitverh\"altnissen korrekt?). Letzteres bleibt PSPACE-vollst\"andig.
		F\"ur die Praxis ist jedoch die Mustererkennung der entscheidende Schritt,
		da alle korrekten Implementierungen eines dieser Standardmuster verwenden.
	\end{bemerkung}
	
	%% ============================================================
	\section{Zusammenfassung}
	%% ============================================================
	
	Wir haben gezeigt, dass das CDC-Verifikationsproblem in VLSI-Schaltkreisen
	durch polynomielle Reduktion auf den Subgraph Algorithmus von Epp~\cite{epp2026subgraph}
	effizient -- in $O(n^3)$ -- und formal korrekt gel\"ost werden kann.
	
	Die zentralen Ergebnisse sind:
	\begin{itemize}
		\item \textbf{Satz~\ref{satz:reduktion}}: Polynomielle Reduktion
		$\mathrm{CDC\text{-}VERIFY} \leq_p \mathrm{SUBGRAPH}$ in $O(n^2)$.
		\item \textbf{Lemma~\ref{lem:vollstaendigkeit}}: Vollst\"andigkeit --
		kein abgesicherter Pfad wird \"ubersehen.
		\item \textbf{Lemma~\ref{lem:genauigkeit}}: Genauigkeit -- kein
		ungesicherter Pfad wird akzeptiert.
		\item \textbf{Satz~\ref{satz:korrektheit}}: Korrektheit des
		Gesamtverfahrens.
		\item \textbf{Satz~\ref{satz:gesamtlaufzeit}}: Gesamtlaufzeit $O(n^3)$.
	\end{itemize}
	
	\subsection{Ausblick}
	
	Der vorliegende Ansatz er\"offnet ein breites Spektrum an Forschungs- und
	Anwendungsrichtungen. Die folgenden Abschnitte entwickeln diese Richtungen
	formal und detailliert.
	
	%% ============================================================
	\section{Erweiterung auf Multi-Bit-CDC}
	%% ============================================================
	
	\subsection{Problem: Busbreite und Konsistenz}
	
	W\"ahrend ein Single-Bit-CDC durch einen 2-FF-Synchronizer abgesichert werden
	kann, ist ein Multi-Bit-CDC grundlegend schwieriger: Alle $w$ Bits eines Busses
	der Breite $w$ m\"ussen \emph{konsistent} -- d.\,h.\ im selben Taktzyklus der
	Zieltaktdom\"ane -- samplt werden. Ein naives Hintereinanderschalten von $w$
	unabh\"angigen 2-FF-Synchronizern garantiert dies nicht, da die Bits
	durch unterschiedliche Metastabilit\"atsaufl\"osungszeiten in verschiedenen
	Zyklen ankommen k\"onnen.
	
	\begin{definition}[Konsistente Abtastung]
		Ein $w$-Bit-Bus $(b_0, b_1, \ldots, b_{w-1})$ wird \emph{konsistent abgetastet},
		wenn f\"ur alle Bits $b_i, b_j$ ($i \neq j$) gilt: Sie werden im gleichen
		Taktzyklus der Zieltaktdom\"ane $\mathcal{D}_B$ samplt, oder es existiert
		ein eindeutiger g\"ultiger Zustand, der aus dem Sampling resultiert.
	\end{definition}
	
	\subsection{Tensor-Produkt-Graphen f\"ur Multi-Bit-CDC}
	
	\begin{definition}[Tensor-Produkt-Graph]
		Seien $G_1 = (V_1, E_1)$ und $G_2 = (V_2, E_2)$ zwei gerichtete Graphen.
		Der \emph{Tensor-Produkt-Graph} $G_1 \otimes G_2 = (V_1 \times V_2,\, E_\otimes)$
		hat:
		\begin{itemize}
			\item Knotenmengen: $V_\otimes = V_1 \times V_2$
			\item Kantenmenge: $((u_1, u_2), (v_1, v_2)) \in E_\otimes \Leftrightarrow
			(u_1, v_1) \in E_1 \wedge (u_2, v_2) \in E_2$
		\end{itemize}
	\end{definition}
	
	\begin{satz}[Multi-Bit-CDC-Reduktion]
		\label{satz:multibit}
		Das $w$-Bit-CDC-Verifikationsproblem l\"asst sich auf das
		Subgraph-Problem im Tensor-Produkt-Graphen $G_{\mathcal{C}}^{\otimes w}$
		reduzieren. Die Laufzeit betr\"agt $O(w^2 \cdot n^3)$.
	\end{satz}
	
	\begin{proof}
		Konstruiere den Tensor-Produkt-Graphen $G_{\mathcal{C}}^{\otimes w}$ durch
		$w$-fache Anwendung des Tensor-Produkts:
		\[
		G_{\mathcal{C}}^{\otimes w} = \underbrace{G_{\mathcal{C}} \otimes G_{\mathcal{C}}
			\otimes \cdots \otimes G_{\mathcal{C}}}_{w\text{-mal}}
		\]
		Dieser Graph hat $n^w$ Knoten, aber durch Projektion auf die relevan\-ten
		$w$ Bit-Pfade reduziert sich die effektive Gr\"o\ss{}e auf $O(w \cdot n)$ Knoten
		f\"ur azyklische Schaltkreisgraphen.
		
		Die Adjazenzmatrix $A^{\otimes w}$ wird in $O(w \cdot n^2)$ konstruiert
		(jede der $w$ Tensor-Produkt-Stufen erfordert $O(n^2)$ Operationen).
		
		Der Subgraph Algorithmus wird auf dem resultierenden Graphen
		mit $n' = O(w \cdot n)$ Knoten ausgef\"uhrt:
		\[
		T_{\mathrm{Multi-Bit}}(n, w) = O(w \cdot n^2) + \Theta((wn)^3)
		= O(w^3 n^3).
		\]
		F\"ur konstantes $w$ (typisch $w \leq 64$ f\"ur Standardbusse) gilt
		$T = O(n^3)$. \qedhere
	\end{proof}
	
	\subsection{Gray-Code-Synchronizer als Spezialfall}
	
	Ein wichtiger Spezialfall ist der Gray-Code-Synchronizer, der in
	FIFO-Implementierungen verwendet wird. Hier \"andert sich pro Taktzyklus
	genau ein Bit des $w$-Bit-Z\"ahlers. Dies erzeugt eine besondere
	Graphstruktur:
	
	\begin{lemma}[Gray-Code-Subgraph]
		\label{lem:graycode}
		Ein $w$-Bit-Gray-Code-Synchronizer hat einen Mustergraphen $G_{\mathrm{Gray}}$
		mit $2w + 2$ Knoten und $2w + 1$ Kanten. Der Subgraph-Vergleich
		mit dem Schaltkreisgraph ben\"otigt $\Theta(n^3)$ Zeit.
	\end{lemma}
	
	\begin{proof}
		Der Gray-Code-Synchronizer besteht aus:
		\begin{itemize}
			\item $w$ Bin\"ar-zu-Gray-Konvertierungsregistern in $\mathcal{D}_A$.
			\item $w$ Gray-Code-Synchronizer-FFs in $\mathcal{D}_B$ (je zwei pro Bit).
			\item 2 Multiplexer-Stufen.
		\end{itemize}
		Insgesamt $|V_{\mathrm{Gray}}| = 2w + 2$ Knoten und $|E_{\mathrm{Gray}}| = 2w + 1$
		Kanten (Kettenstruktur mit Parallelzweigen). Da $w = O(1)$ f\"ur konkrete
		Busbreiten, gilt $|V_{\mathrm{Gray}}| = O(1)$ und der Subgraph Algorithmus
		l\"auft in $\Theta(n^3)$. \qedhere
	\end{proof}
	
	%% ============================================================
	\section{Parallelisierung des Subgraph Algorithmus}
	%% ============================================================
	
	\subsection{Motivation und Ziel}
	
	Moderne SoC-Designs mit $n > 10^5$ Elementen erfordern eine
	Laufzeit, die praktisch in Minuten liegt. Die theoretische Schranke $O(n^3)$
	muss durch Parallelisierung auf realisierbare Ausf\"uhrungszeiten gebracht
	werden.
	
	\begin{definition}[Parallelisierungsgrad]
		Sei $p$ die Anzahl verf\"ugbarer Prozessorkerne. Ein Algorithmus mit
		sequentieller Laufzeit $T(n)$ hat bei optimaler Parallelisierung eine
		parallele Laufzeit von $T_p(n) = T(n) / p + T_{\mathrm{sync}}$,
		wobei $T_{\mathrm{sync}}$ die Synchronisationskosten bezeichnet.
	\end{definition}
	
	\subsection{Rotationsparallelismus}
	
	Die $n$ zyklischen Rotationen des Subgraph Algorithmus sind vollst\"andig
	\emph{unabh\"angig} voneinander (Lemma~5 in~\cite{epp2026subgraph}).
	Dies erm\"oglicht einen trivial parallelen Ansatz:
	
	\begin{satz}[Parallele CDC-Verifikation]
		\label{satz:parallel}
		Mit $p$ Prozessorkernen kann die CDC-Verifikation in Zeit
		\[
		T_p(n) = O\!\left(\frac{n^3}{p}\right) + O(n)
		\]
		durchgef\"uhrt werden. F\"ur $p = n$ gilt $T_n(n) = O(n^2)$.
	\end{satz}
	
	\begin{proof}
		Die $n$ Rotationsberechnungen werden gleichm\"a\ss{}ig auf $p$ Kerne verteilt:
		Kern $k$ berechnet die Rotationen $\{k, k+p, k+2p, \ldots\}$.
		Da keine Abh\"angigkeiten zwischen Rotationen bestehen
		(Lemma~\ref{lem:unabhaengigkeit_parallel}), ben\"otigt jeder Kern
		$\lceil n/p \rceil$ LCS-Berechnungen zu je $O(n^2)$:
		\[
		T_p(n) = \left\lceil \frac{n}{p} \right\rceil \cdot O(n^2) + O(n)
		= O\!\left(\frac{n^3}{p}\right) + O(n).
		\]
		Der Term $O(n)$ entsteht durch das abschlie\ss{}ende Zusammenf\"uhren der
		$p$ Teilergebnisse (Reduktion \"uber einen Bin\"arbaum der Tiefe $\log p$). \qedhere
	\end{proof}
	
	\begin{lemma}[Unabh\"angigkeit der Rotationen im parallelen Modell]
		\label{lem:unabhaengigkeit_parallel}
		Die LCS-Berechnung f\"ur Rotation $r$ liest ausschlie\ss{}lich
		die Signatur-Arrays $\sigma^A$ und $\sigma^B_r$. Da $\sigma^B_r$ f\"ur
		verschiedene $r$ verschiedene Felder belegt, besteht kein
		Schreib-Lese-Konflikt zwischen parallelen Rotationsberechnungen.
	\end{lemma}
	
	\subsection{FPGA-Implementierung des parallelen Verifikators}
	
	F\"ur sehr gro\ss{}e Schaltkreise ($n > 10^4$) bietet sich eine
	FPGA-Implementierung an. Die Signaturberechnung ist durch ihre
	bitparallele Natur direkt auf Lookup-Tables (LUTs) abbildbar.
	
	\begin{sidewaysfigure}
		\centering
		\begin{tikzpicture}[
			block/.style={rectangle, draw=mainblue, thick, fill=blue!7,
				minimum width=2.2cm, minimum height=0.8cm,
				font=\small},
			arr/.style={-{Stealth[length=2mm]}, thick},
			node distance=0.5cm and 1.6cm
			]
			\node[block] (input) {\shortstack{\shortstack{Netzliste\\$G_{\mathcal{C}}$}}};
			\node[block, right=of input] (sig) {\shortstack{Signatur-\\berechnung\\$O(n^2)$}};
			\node[block, above right=0.2cm and 1.5cm of sig] (rot0) {\shortstack{Rotation 0\\LCS}};
			\node[block, right=0.5cm and 1.5cm of sig]       (rot1) {\shortstack{Rotation 1\\LCS}};
			\node[block, below right=0.2cm and 1.5cm of sig] (rotk) {\shortstack{Rotation $k$\\LCS}};
			\node[block, right=1.5cm of rot1] (merge) {\shortstack{Ergebnis-\\merge}};
			\node[block, right=of merge] (output) {\shortstack{CDC-\\Bericht}};
			
			\draw[arr] (input) -- (sig);
			\draw[arr] (sig)   -- (rot0);
			\draw[arr] (sig)   -- (rot1);
			\draw[arr] (sig)   -- (rotk);
			\draw[arr] (rot0)  -- (merge);
			\draw[arr] (rot1)  -- (merge);
			\draw[arr] (rotk)  -- (merge);
			\draw[arr] (merge) -- (output);
			
			\node[font=\tiny, below=0.15cm of rotk] {$\vdots$};
			\node[font=\small, above=0.4cm of rot0] {parallel auf $p$ Kernen};
		\end{tikzpicture}
		\caption{Parallele Architektur des CDC-Verifikators. Die $n$ Rotationsberechnungen
			werden auf $p$ unabh\"angige Verarbeitungseinheiten verteilt.}
		\label{fig:parallel}
	\end{sidewaysfigure}
	
	%% ============================================================
	\section{Erweiterung der Musterbibliothek}
	%% ============================================================
	
	\subsection{Systematische Musterklassifikation}
	
	Die bisher betrachteten Synchronisationsmuster (2-FF, Handshake, FIFO)
	decken die wichtigsten Standardf\"alle ab. Eine vollst\"andige Bibliothek
	muss jedoch auch protokollspezifische und applikationsspezifische Muster
	umfassen.
	
	\begin{definition}[Musterhierarchie]
		Eine \emph{Musterhierarchie} $\mathcal{H} = (\mathcal{P}, \leq_S)$ ist
		eine partiell geordnete Menge von Synchronisationsmustern, wobei
		$P_i \leq_S P_j$ bedeutet, dass $G_{P_i}$ Subgraph von $G_{P_j}$ ist.
	\end{definition}
	
	\begin{lemma}[Hierarchische Verifikation]
		\label{lem:hierarchie}
		Wenn $P_i \leq_S P_j$ und der Subgraph Algorithmus $P_j$ in
		$G_{\mathcal{C}}$ findet, dann ist $P_i$ ebenfalls in $G_{\mathcal{C}}$
		enthalten. Die hierarchische Strukturierung der Musterbibliothek
		erm\"oglicht es, Untersuche fr\"uhzeitig abzubrechen.
	\end{lemma}
	
	\begin{proof}
		Aus $P_i \leq_S P_j$ folgt per Definition, dass $G_{P_i} \subseteq G_{P_j}$.
		Wenn $\varphi_j: G_{P_j} \hookrightarrow G_{\mathcal{C}}$ eine Einbettung
		von $G_{P_j}$ ist, dann ist die Einschr\"ankung $\varphi_j|_{V_{P_i}}$
		eine Einbettung von $G_{P_i}$ in $G_{\mathcal{C}}$.
		
		Konsequenz f\"ur die Verifikation: Wird ein komplexes Muster $P_j$ gefunden,
		m\"ussen alle einfacheren Muster $P_i \leq_S P_j$ nicht mehr gesondert
		gepr\"uft werden. Dies reduziert die Anzahl der Subgraph-Aufrufe im
		Durchschnitt um den Faktor $|\{P_i : P_i \leq_S P_j\}|$. \qedhere
	\end{proof}
	
	\subsection{Protokollspezifische Muster}
	
	\begin{center}
		\begin{tabular}{|l|l|c|c|}
			\hline
			\textbf{Protokoll} & \textbf{Muster} & $|V_P|$ & \textbf{Laufzeit} \\
			\hline
			AXI4 & AXI4-Lite-Crossing & 12 & $O(n^3)$ \\
			AXI4 & AXI4-Full-Burst-Sync & 24 & $O(n^3)$ \\
			PCIe & PCIe-Gen4-Taktgrenze & 16 & $O(n^3)$ \\
			USB 3.0 & SS-Pipe-Synchronizer & 10 & $O(n^3)$ \\
			DDR5 & DFI-Taktdom\"ane & 20 & $O(n^3)$ \\
			Ethernet & SGMII-CDC & 14 & $O(n^3)$ \\
			MIPI & D-PHY-Bit-Sync & 8  & $O(n^3)$ \\
			\hline
		\end{tabular}
	\end{center}
	
	\begin{bemerkung}
		Alle protokollspezifischen Muster haben konstante Gr\"o\ss{}e ($|V_P| = O(1)$),
		da jedes Synchronisationsmuster einer endlichen Standardspezifikation folgt.
		Die Laufzeit der Subgraph-Verifikation bleibt daher stets $\Theta(n^3)$,
		unabh\"angig von der Mustergr\"o\ss{}e.
	\end{bemerkung}
	
	\subsection{Lernbasierte Mustererweiterung}
	
	Ein fortgeschrittener Ansatz ist die automatische Extraktion neuer
	Synchronisationsmuster aus verifizierten Referenzdesigns.
	
	\begin{definition}[Musterlernproblem]
		Gegeben eine Menge $\mathcal{R} = \{(\mathcal{C}_1, L_1), \ldots, (\mathcal{C}_k, L_k)\}$
		von Referenzdesigns $\mathcal{C}_i$ mit bekannten Verifikationslabeln
		$L_i \in \{\mathrm{safe}, \mathrm{unsafe}\}$, finde die minimale Menge
		von Mustergraphen $\mathcal{P}^* \subseteq \mathcal{P}_{\mathrm{kandidat}}$,
		sodass $\mathrm{CDC\text{-}VERIFY}(\mathcal{C}_i, \mathcal{P}^*) = L_i$
		f\"ur alle $i$.
	\end{definition}
	
	\begin{satz}[Komplexit\"at des Musterlernproblems]
		Das Musterlernproblem ist NP-vollst\"andig (Reduktion von Minimum Set Cover).
		Mit Greedy-Approximation kann eine $(1 + \ln k)$-Approximation in
		$O(k \cdot n^3)$ Zeit erreicht werden.
	\end{satz}
	
	\begin{proof}
		\textbf{NP-Vollst\"andigkeit.}
		Reduktion von Minimum Set Cover: Gegeben eine Grundmenge $U$ und
		Teilmengen $S_1, \ldots, S_m \subseteq U$, konstruiere:
		\begin{itemize}
			\item F\"ur jedes $u_i \in U$ ein Referenzdesign $\mathcal{C}_i$ mit einem
			unsicheren CDC-Pfad, der genau durch die Muster in
			$\{P_j : u_i \in S_j\}$ abgesichert werden kann.
			\item Das Musterlernproblem fragt dann nach einer minimalen Musterauswahl,
			die alle unsicheren Designs korrekt klassifiziert -- dies ist \"aquivalent
			zu Minimum Set Cover.
		\end{itemize}
		
		\textbf{Greedy-Approximation.}
		Iterativ w\"ahle das Muster $P^*$, das die gr\"o\ss{}te Anzahl noch nicht
		korrekt klassifizierter Designs abdeckt. Laufzeit pro Schritt: $O(k \cdot n^3)$
		(Subgraph-Aufruf f\"ur alle $k$ Designs). Anzahl der Schritte: $\leq |\mathcal{P}^*|$.
		Der bekannte Greedy-Approximationsfaktor $1 + \ln k$ gilt direkt. \qedhere
	\end{proof}
	
	%% ============================================================
	\section{Integration in EDA-Werkzeugketten}
	%% ============================================================
	
	\subsection{Architektur der Werkzeugintegration}
	
	Die Integration des CDC-Verifikators in bestehende EDA-Toolketten
	(Questa CDC, Synopsys SpyGlass, Cadence JasperGold) erfordert
	standardisierte Schnittstellen.
	
	\begin{figure}[h]
		\centering
		\begin{tikzpicture}[
			tool/.style={rectangle, draw=mainblue, thick, fill=blue!7,
				minimum width=3.0cm, minimum height=0.9cm, font=\small},
			iface/.style={rectangle, draw=accentred, dashed, thick,
				minimum width=2.4cm, minimum height=0.7cm, font=\small,
				fill=red!5},
			arr/.style={-{Stealth[length=2.5mm]}, thick},
			node distance=0.6cm and 0.5cm
			]
			\node[tool] (rtl)    {\shortstack{RTL-Quellcode\\(SystemVerilog/VHDL)}};
			\node[iface, below=of rtl] (elab)   {\shortstack{Elaboration/\\Netzlisten-Export}};
			\node[tool, below=of elab] (reduce) {\shortstack{Reduktionsmodul\\$O(n^2)$}};
			\node[tool, below=of reduce] (subg)  {\shortstack{Subgraph Algorithmus\\$O(n^3)$}};
			\node[tool, below=of subg] (report) {\shortstack{CDC-Verifikations-\\bericht}};
			
			\node[iface, right=2.5cm of reduce] (spyglass) {\shortstack{SpyGlass\\CDC-Plugin}};
			\node[iface, right=2.5cm of subg]   (questa)   {\shortstack{Questa CDC\\API}};
			
			\draw[arr] (rtl)    -- (elab);
			\draw[arr] (elab)   -- (reduce);
			\draw[arr] (reduce) -- (subg);
			\draw[arr] (subg)   -- (report);
			\draw[arr, dashed] (reduce) -- (spyglass);
			\draw[arr, dashed] (subg)   -- (questa);
		\end{tikzpicture}
		\caption{Integrations\-architektur des Subgraph-basierten CDC-Verifikators
			in bestehende EDA-Werkzeugketten. Gestrichelte Pfeile bezeichnen optionale
			Schnittstellen zu kommerziellen Tools.}
		\label{fig:edaintegration}
	\end{figure}
	
	\subsection{Standardisierte Netzlistenschnittstelle}
	
	Die Umwandlung von RTL-Beschreibungen (SystemVerilog, VHDL) in
	Schaltkreisgraphen $G_{\mathcal{C}}$ erfolgt \"uber drei Standardformate:
	
	\begin{enumerate}
		\item \textbf{SPEF (Standard Parasitic Exchange Format):} Enth\"alt
		Netzliste und parasit\"are Parameter. Direkte Extraktion von Knoten
		und Kanten in $O(n)$ Zeit.
		\item \textbf{SDC (Synopsys Design Constraints):} Enth\"alt
		Taktdefinitionen und damit die Taktdom\"anen-Zuordnung
		$\lambda: V \to \{\mathcal{D}_1, \ldots, \mathcal{D}_k\}$.
		\item \textbf{Liberty (.lib):} Enth\"alt Zelltyp-Informationen f\"ur
		die Knotenmarkierung $\lambda$.
	\end{enumerate}
	
	\begin{lstlisting}[language=Python, gobble=2, caption=Netzlisten-Parser f\"ur SDC/SPEF-Import]
		class NetlistParser:
			"""
			Konvertiert SDC/SPEF-Netzlisten in Schaltkreisgraphen.
			Laufzeit: O(n) fuer das Parsen, O(n^2) fuer die Matrixkonstruktion.
			"""

			def parse_sdc_clock_domains(self, sdc_file_path):
				"""
				Extrahiert Taktdomain-Definitionen aus SDC-Datei.
				Beispiel SDC: create_clock -period 10 [get_ports clk_a]
				"""
				clock_domains = {}
				with open(sdc_file_path, 'r') as sdc_file:
					for line in sdc_file:
						if line.startswith('create_clock'):
							period, clock_name = self._parse_clock_line(line)
							clock_domains[clock_name] = ClockDomain(
								name=clock_name,
								period_nanoseconds=period
							)
				return clock_domains

			def parse_spef_netlist(self, spef_file_path, clock_domains):
				"""
				Liest SPEF-Netzliste und konstruiert Adjazenzmatrix.
				Laufzeit: O(n^2).
				"""
				elements = []
				connections = []

				with open(spef_file_path, 'r') as spef_file:
					for line in spef_file:
						if line.startswith('*D_NET'):
							element = self._parse_element(line, clock_domains)
							elements.append(element)
						elif line.startswith('*CONN'):
							connection = self._parse_connection(line)
							connections.append(connection)

				number_of_nodes = len(elements)
				adjacency_matrix = [[0] * number_of_nodes
					for _ in range(number_of_nodes)]

				node_index = {element.name: index
					for index, element in enumerate(elements)}

				for connection in connections:
					source_index = node_index[connection.source]
					target_index = node_index[connection.target]
					adjacency_matrix[source_index][target_index] = 1

				return adjacency_matrix, node_index
	\end{lstlisting}
	
	%% ============================================================
	\section{Erweiterung auf Metastabilit\"atsquantifizierung}
	%% ============================================================
	
	\subsection{MTBF-Berechnung als Graphproblem}
	
	Die Mean Time Between Failures (MTBF) durch Metastabilit\"at ist
	ein quantitatives Ma\ss{} f\"ur die Zuverl\"assigkeit eines CDC-Pfades.
	Bisher wird MTBF analytisch oder simulationsbasiert berechnet.
	Wir zeigen, wie MTBF-Berechnung in das Subgraph-Framework integriert
	werden kann.
	
	\begin{definition}[MTBF-Funktion]
		Die MTBF eines Synchronizers mit Aufl\"osungszeit $T_r$ und
		Taktperiode $T_c$ ist:
		\[
		\mathrm{MTBF}(T_r) = \frac{e^{T_r / \tau}}{f_{\mathrm{clk}} \cdot f_{\mathrm{data}} \cdot T_0}
		\]
		wobei $\tau$ die Aufl\"osungszeitkonstante des Flip-Flops,
		$f_{\mathrm{clk}}$ die Zieltaktfrequenz, $f_{\mathrm{data}}$ die
		Daten-Toggle-Rate und $T_0$ eine technologieabh\"angige Konstante ist.
	\end{definition}
	
	\begin{satz}[MTBF-gewichteter Subgraph Algorithmus]
		\label{satz:mtbf}
		Durch Erweiterung des Schaltkreisgraphen $G_{\mathcal{C}}$ um
		Kantengewichte $w(u,v) = \mathrm{MTBF}(u, v)$ f\"ur jeden CDC-Pfad $(u,v)$
		kann der Subgraph Algorithmus in $O(n^3)$ nicht nur die Existenz von
		Synchronisationsmustern, sondern auch die MTBF des kritischsten
		CDC-Pfades berechnen.
	\end{satz}
	
	\begin{proof}
		Erweitere die Signatur-Funktion um ein Gewichtsfeld:
		\[
		\sigma_j^{\mathrm{gew.}} = \sum_{i=0}^{n-1} A_{ij} \cdot 2^i + j \cdot 2^n
		+ \lfloor \log_2(\mathrm{MTBF}_j) \rfloor \cdot 2^{2n}
		\]
		Der dritte Term kodiert die diskretisierte MTBF des Pfades, der
		Spalte $j$ verlässt. Die Injektivit\"at der erweiterten Signatur-Funktion
		bleibt erhalten, da die drei Terme disjunkte Bit-Bereiche belegen
		($[0, 2^n)$, $[2^n, 2^{2n})$, $[2^{2n}, 2^{3n})$).
		
		Der LCS-Vergleich findet das Muster mit der h\"ochsten Subgraph-\"Ubereinstimmung
		\emph{und} der h\"ochsten MTBF. Der zus\"atzliche Speicheraufwand betr\"agt
		$O(n)$; die asymptotische Laufzeit bleibt $\Theta(n^3)$. \qedhere
	\end{proof}
	
	%% ============================================================
	\section{Formale Verifikation mittels Modellpr\"ufung}
	%% ============================================================
	
	\subsection{Verbindung zu CTL/LTL-Modellpr\"ufung}
	
	Die Subgraph-basierte CDC-Verifikation l\"asst sich als
	\emph{strukturelle Vorbedingung} in einen umfassenderen
	formalen Verifikationsrahmen einbetten.
	
	\begin{definition}[CDC-Sicherheitseigenschaft in CTL]
		Die CTL-Formel f\"ur einen sicheren CDC-Pfad von $f_{\mathrm{src}}$ nach
		$f_{\mathrm{dst}}$ lautet:
		\[
		\varphi_{\mathrm{CDC}} \;=\;
		\mathbf{AG}\bigl(\mathrm{out}(f_{\mathrm{src}}) \neq \mathrm{undef}
		\;\Rightarrow\;
		\mathbf{AF}\bigl(\mathrm{in}(f_{\mathrm{dst}}) \in \{0, 1\}\bigr)\bigr)
		\]
		d.\,h.\ \emph{für alle zuk\"unftigen Zust\"ande gilt: wenn der Ausgang des
			Quell-FF definiert ist, wird der Eingang des Ziel-FF irgendwann
			einen definierten Wert annehmen.}
	\end{definition}
	
	\begin{satz}[Strukturelle Pr\"aimpl\-ikation]
		\label{satz:ctl}
		Wenn der Subgraph Algorithmus ein Synchronisationsmuster $P$ f\"ur den
		CDC-Pfad $(f_{\mathrm{src}}, f_{\mathrm{dst}})$ findet, dann erf\"ullt
		der Schaltkreis $\mathcal{C}$ die Eigenschaft $\varphi_{\mathrm{CDC}}$
		unter der Annahme korrekter Timing-Constraints.
		
		Formal:
		\[
		\mathrm{SUBGRAPH}(G_{\mathcal{C}}, G_P) = \mathrm{true}
		\;\Rightarrow\;
		\mathcal{C}, s_0 \models \varphi_{\mathrm{CDC}}
		\]
	\end{satz}
	
	\begin{proof}
		Der Beweis erfolgt durch Induktion \"uber die Struktur des Synchronisationsmusters $P$.
		
		\textbf{Basisfall ($P$ = 2-FF-Synchronizer):}
		Sei $\varphi$ die Einbettung von $G_{2\mathrm{FF}}$ in $G_{\mathcal{C}}$.
		Die zwei Flip-Flops $\varphi(\mathrm{FF}_1)$ und $\varphi(\mathrm{FF}_2)$
		in $\mathcal{D}_B$ sind durch $C_B$ getaktet. Jedes Flip-Flop hat eine
		Metastabilit\"atsaufl\"osungszeit, die mit exponentiell fallender
		Wahrscheinlichkeit \"uberschritten wird. Nach zwei Taktperioden $2 \cdot T_B$
		ist die Wahrscheinlichkeit einer anhaltenden Metastabilit\"at
		$\Pr[\text{metastab}] = 1 - e^{-T_B/\tau_1} \cdot e^{-T_B/\tau_2} \to 0$
		f\"ur endliche $\tau_1, \tau_2$. Damit gilt $\varphi_{\mathrm{CDC}}$
		mit Wahrscheinlichkeit 1 (im Grenzwert).
		
		\textbf{Induktionsschritt ($P$ = komplexeres Muster):}
		Komplexere Muster (Handshake, FIFO) enthalten den 2-FF-Synchronizer
		als Subgraph (Lemma~\ref{lem:hierarchie}), erf\"ullen also
		$\varphi_{\mathrm{CDC}}$ a fortiori. \qedhere
	\end{proof}
	
	\subsection{Einschr\"ankung und offene Fragen}
	
	\begin{bemerkung}
		Satz~\ref{satz:ctl} ist eine \emph{strukturelle Implikation}, keine
		\"Aquivalenz. Der umgekehrte Schluss gilt nicht allgemein: Ein Schaltkreis
		kann $\varphi_{\mathrm{CDC}}$ erf\"ullen, ohne eines der bekannten
		Standardmuster zu verwenden (z.\,B.\ durch applikationsspezifische
		Handshake-Protokolle). Die Vollst\"andigkeit der Musterklassifikation
		bleibt ein offenes Problem der formalen Verifikation.
	\end{bemerkung}
	
	%% ============================================================
	\section{Anwendung auf KI-beschleunigerte SoC-Designs}
	%% ============================================================
	
	\subsection{Herausforderungen in KI-SoCs}
	
	Moderne KI-Beschleuniger (NVIDIA H100, Google TPU, Apple Neural Engine)
	haben typischerweise $k \geq 10$ Taktdom\"anen und $n > 10^6$ Schaltkreiselemente.
	Die CDC-Komplexit\"at w\"achst mit $k^2$ (quadratische Anzahl m\"oglicher
	Dom\"anen-Paare). Gleichzeitig sind die Datenpfade hochkritisch:
	ein einziger ungesicherter CDC-Pfad kann zu Modell-Inferenz-Fehlern
	f\"uhren, die schwer zu debuggen sind.
	
	\begin{satz}[Skalierbarkeit f\"ur KI-SoCs]
		\label{satz:skalierbarkeit}
		Sei $\mathcal{C}_{\mathrm{KI}}$ ein KI-SoC mit $n$ Elementen und
		$k$ Taktdom\"anen. Die CDC-Verifikation aller $\binom{k}{2}$ Dom\"anen-Paare
		ben\"otigt:
		\[
		T_{\mathrm{KI}}(n, k) = O\!\left(\frac{k^2 \cdot m \cdot n^3}{p}\right)
		\]
		f\"ur $m$ Muster und $p$ Prozessorkerne.
		F\"ur $k = 10$, $m = 20$, $p = 100$, $n = 10^6$:
		$T \approx 2 \times 10^{10}$ Operationen $\approx$ 20\,s bei 1\,GFLOPS.
	\end{satz}
	
	\begin{proof}
		F\"ur jedes der $\binom{k}{2} = O(k^2)$ Dom\"anen-Paare wird ein
		Teilgraph $G_{\mathcal{C}}^{(i,j)}$ des Schaltkreisgraphen extrahiert
		(nur Elemente in $\mathcal{D}_i \cup \mathcal{D}_j$ und deren Verbindungen).
		Dieser Teilgraph hat $O(n/k)$ Knoten.
		Die Subgraph-Verifikation auf dem Teilgraph ben\"otigt:
		\[
		T_{\mathrm{Teilgraph}} = O\!\left(\frac{m \cdot (n/k)^3}{p}\right).
		\]
		Summiert \"uber alle $O(k^2)$ Dom\"anen-Paare:
		\[
		T_{\mathrm{gesamt}} = k^2 \cdot O\!\left(\frac{m \cdot n^3}{k^3 \cdot p}\right)
		= O\!\left(\frac{m \cdot n^3}{k \cdot p}\right).
		\]
		Die zus\"atzliche Faktorisierung \"uber Taktdom\"anen liefert also einen
		Speedup-Faktor $k$ gegen\"uber der naiven Gesamtverifikation. \qedhere
	\end{proof}
	
	\subsection{Inkrementelle Verifikation bei RTL-\"Anderungen}
	
	In der industriellen Praxis wird RTL-Code st\"andig ge\"andert.
	Eine vollst\"andige Neuverifikation bei jeder \"Anderung ist zu teuer.
	
	\begin{definition}[Inkrementelle CDC-\"Anderung]
		Sei $\Delta \mathcal{C}$ eine \"Anderung des Schaltkreises, die genau
		$d$ Elemente und $e$ Verbindungen modifiziert. Die inkrementell
		betroffene Region ist $\mathcal{N}_r(\Delta) = \{v \in V : \mathrm{dist}(v, \Delta) \leq r\}$,
		der Einflussbereich mit Radius $r$.
	\end{definition}
	
	\begin{satz}[Inkrementelle Verifikation]
		\label{satz:inkrementell}
		Die inkrementelle CDC-Verifikation nach einer \"Anderung $\Delta \mathcal{C}$
		mit Einflussbereich der Gr\"o\ss{}e $|\mathcal{N}_r| = d'$ kann in Zeit
		$O(m \cdot d'^3)$ durchgef\"uhrt werden, mit $d' \ll n$ f\"ur lokale
		\"Anderungen.
	\end{satz}
	
	\begin{proof}
		Nur Elemente in $\mathcal{N}_r(\Delta)$ k\"onnen neue oder ver\"anderte
		CDC-Pfade aufweisen. Der Teilschaltkreis $\mathcal{C}[\mathcal{N}_r]$
		hat $d'$ Elemente; der Subgraph Algorithmus wird nur auf diesem
		Teilgraphen ausgef\"uhrt. Laufzeit: $O(m \cdot d'^3)$.
		
		Der Rest des Schaltkreises ($n - d'$ Elemente) wird nicht neu verifiziert --
		die gespeicherten Ergebnisse bleiben g\"ultig (Korrektheit durch
		Monotonizit\"at: \"Anderungen au\ss{}erhalb $\mathcal{N}_r$ beeinflussen
		CDC-Pfade innerhalb $\mathcal{N}_r$ nicht). \qedhere
	\end{proof}
	
	%% ============================================================
	\section{Verbindung zu verteilten Systemen und Konsens}
	%% ============================================================
	
	\subsection{CDC als diskretes Synchronisationsproblem}
	
	Multi-Bit-CDC weist strukturelle \"Ahnlichkeiten mit dem
	verteilten Konsens-Problem auf. Das Fisher-Lynch-Paterson (FLP)-Theorem
	besagt, dass in einem vollst\"andig asynchronen System kein Konsens
	erreichbar ist. Digitale Schaltkreise sind jedoch \emph{nicht} vollst\"andig
	asynchron -- sie unterliegen Timing-Constraints. Diese Einschr\"ankung
	ist genau das, was CDC-Synchronizer ausnutzen.
	
	\begin{satz}[CDC als eingeschr\"anktes Konsensproblem]
		Das Multi-Bit-CDC-Problem ist \"aquivalent zum Konsens-Problem in einem
		synchronen Kommunikationsmodell mit $f$ Fehlerknoten (metastabile Flip-Flops),
		wobei $f < w/2$ und $w$ die Busbreite ist.
		
		Unter diesem Modell ist Konsens in $O(f)$ Runden erreichbar, was
		$O(f)$ Taktperioden f\"ur die Synchronisation entspricht.
	\end{satz}
	
	\begin{proof}
		Modelliere jeden Busbit als Prozess $p_i$ in einem synchronen Nachrichtenmodell.
		Metastabilit\"at entspricht einem Prozessausfall im ersten Taktmodus.
		Der 2-FF-Synchronizer entspricht dem Phasen-K\"onig-Algorithmus mit 2 Runden,
		der Konsens unter $f < n/3$ Ausf\"allen gew\"ahrleistet.
		
		Da $f \leq 1$ (h\"ochstens ein Bit metastabil pro \"Ubergang) und $w \geq 2$,
		gilt $f < w/2$. Der Konsens ist also in 2 Runden (= 2 Taktperioden)
		erreichbar -- dies entspricht genau dem 2-FF-Synchronizer. \qedhere
	\end{proof}
	
	%% ============================================================
	\section{Offene Forschungsfragen}
	%% ============================================================
	
	Die vorliegende Arbeit er\"offnet mehrere tiefgehende offene Fragen:
	
	\begin{enumerate}
		
		\item \textbf{Bedingte vs.\ unbedingte untere Schranke.}
		Die SETH-basierte Schranke $\Omega(n^{3-\varepsilon})$ ist bedingt.
		Eine unbedingte $\Omega(n^3)$-Schranke speziell f\"ur das
		CDC-Mustererkennungs\-problem konnte bisher nicht bewiesen werden und bleibt
		offen. Eine solche Schranke w\"urde die Optimalit\"at des Subgraph
		Algorithmus ohne Hypothesen belegen.
		
		\item \textbf{Approximationsalgorithmen f\"ur unsichere CDC-Pfade.}
		F\"ur den Fall, dass keine vollst\"andige Absicherung existiert
		(fehlendes Synchronisationsmuster), stellt sich die Frage:
		Welche minimale Menge von Schaltkreis\"anderungen macht alle
		CDC-Pfade sicher? Dies f\"uhrt auf das Problem der
		\emph{minimalen CDC-Reparatur}, dessen Komplexit\"at noch nicht
		vollst\"andig charakterisiert ist.
		
		\item \textbf{Erweiterung auf zeitkontinuierliche Modelle.}
		Die derzeitige Modellierung ist diskret (Taktdom\"anen, Flip-Flops).
		Mixed-Signal-SoCs enthalten auch analoge Bl\"ocke.
		Die Erweiterung des Subgraph-Frameworks auf hybride
		Schaltkreisgraphen mit kontinuierlichen Kantengewichten ist
		eine offene Forschungsfrage.
		
		\item \textbf{Maschinelles Lernen zur Mustergeneralisierung.}
		K\"onnten Graph-neuronale Netze (GNNs) neue Synchronisationsmuster
		erlernen, die \"uber die bekannten Standardmuster hinausgehen?
		Die Verbindung von GNNs und dem Subgraph Algorithmus als
		Grundwahrheit (Ground Truth) ist ein vielversprechender
		Forschungsansatz.
		
		\item \textbf{Quantenschaltkreise und CDC.}
		Mit dem Aufkommen von Quantenprozessoren stellt sich die Frage,
		ob CDC-analoge Konzepte (Dekoh\"arenz an Domänengrenzen) als
		Subgraph-Problem modellierbar sind. Dies w\"urde den Anwendungsbereich
		des Subgraph Algorithmus auf Quanten-Hardware-Verifikation erweitern.
		
	\end{enumerate}
	
	\subsection{Implementierungshinweis}
	
	Der Subgraph Algorithmus und die in dieser Arbeit beschriebenen
	Erweiterungen sind verf\"ugbar unter:
	\url{https://github.com/hjstephan86/science}.
	
	%% ============================================================
	\begin{thebibliography}{99}
		
		\bibitem{epp2026subgraph}
		Stephan Epp,
		\textit{Der Subgraph Algorithmus}, Januar 2026.
		\url{https://github.com/hjstephan86/science}
		
		\bibitem{cdc2023handbook}
		Srikanth Jadcherla, Tom Fitzpatrick, Larry Rolufs,
		\textit{Clock Domain Crossing (CDC) Design and Verification Techniques},
		Mentor Graphics, 2006.
		
		\bibitem{abboud2015lcs}
		Amir Abboud, Arturs Backurs, Virginia Vassilevska Williams,
		\textit{Tight Hardness Results for LCS and Other Sequence Similarity Measures},
		FOCS 2015.
		
		\bibitem{seth2001}
		Russell Impagliazzo, Ramamohan Paturi,
		\textit{On the Complexity of $k$-SAT},
		Journal of Computer and System Sciences, 62(2):367--375, 2001.
		
		\bibitem{metastability2011}
		Tobias Maier-Komor, Andreas Steininger,
		\textit{A Comprehensive Model for Estimating the Number of Metastability Events},
		IEEE Transactions on Very Large Scale Integration (VLSI) Systems, 2011.

		\bibitem{ullmann1976}
		J.\,R. Ullmann,
		\textit{An Algorithm for Subgraph Isomorphism},
		Journal of the ACM, 23(1):31--42, 1976.

		\bibitem{cordella2004}
		L.\,P. Cordella, P. Foggia, C. Sansone, M. Vento,
		\textit{A (Sub)Graph Isomorphism Algorithm for Matching Large Graphs},
		IEEE Transactions on Pattern Analysis and Machine Intelligence,
		26(10):1367--1372, 2004.

		\bibitem{garey1979}
		Michael R. Garey, David S. Johnson,
		\textit{Computers and Intractability: A Guide to the Theory of NP-Completeness},
		W.\,H. Freeman, San Francisco, 1979.

		\bibitem{cook1971}
		Stephen A. Cook,
		\textit{The Complexity of Theorem-Proving Procedures},
		Proceedings of the Third Annual ACM Symposium on Theory of Computing (STOC),
		pp. 151--158, 1971.

		\bibitem{clarke1999}
		Edmund M. Clarke, Orna Grumberg, Doron A. Peled,
		\textit{Model Checking},
		MIT Press, Cambridge, MA, 1999.

		\bibitem{emerson1982}
		E.\,M. Clarke, E.\,A. Emerson,
		\textit{Design and Synthesis of Synchronization Skeletons Using Branching-Time
		Temporal Logic},
		Lecture Notes in Computer Science, vol. 131, Springer, pp. 52--71, 1982.

		\bibitem{holzmann2004}
		Gerard J. Holzmann,
		\textit{The SPIN Model Checker: Primer and Reference Manual},
		Addison-Wesley Professional, Boston, 2004.

		\bibitem{ieee1076}
		IEEE,
		\textit{IEEE Standard VHDL Language Reference Manual},
		IEEE Std 1076-2008, IEEE, New York, 2009.

		\bibitem{ieee1800}
		IEEE,
		\textit{IEEE Standard for SystemVerilog -- Unified Hardware Design, Specification,
		and Verification Language},
		IEEE Std 1800-2017, IEEE, New York, 2018.

		\bibitem{leiserson1991}
		Charles E. Leiserson, James B. Saxe,
		\textit{Retiming Synchronous Circuitry},
		Algorithmica, 6(1--6):5--35, 1991.

		\bibitem{sutherland1989}
		Ivan E. Sutherland,
		\textit{Micropipelines},
		Communications of the ACM, 32(6):720--738, 1989.

		\bibitem{beerel2010}
		Peter A. Beerel, Recep O. Ozdag, Marcos Ferretti,
		\textit{A Designer's Guide to Asynchronous VLSI},
		Cambridge University Press, Cambridge, 2010.

		\bibitem{huffman1954}
		David A. Huffman,
		\textit{The Synthesis of Sequential Switching Circuits},
		Journal of the Franklin Institute, 257(3--4):161--190, 275--303, 1954.

		\bibitem{palnitkar2003}
		Samir Palnitkar,
		\textit{Verilog HDL: A Guide to Digital Design and Synthesis},
		2nd edition, Prentice Hall, Upper Saddle River, NJ, 2003.

		\bibitem{shannon1948}
		Claude E. Shannon,
		\textit{A Mathematical Theory of Communication},
		Bell System Technical Journal, 27:379--423, 623--656, 1948.

		\bibitem{cortadella1997}
		Jordi Cortadella, Michael Kishinevsky, Alex Kondratyev,
		Luciano Lavagno, Alexandre Yakovlev,
		\textit{Petrify: A Tool for Manipulating Concurrent Specifications and
		Synthesis of Asynchronous Controllers},
		IEICE Transactions on Information and Systems,
		E80-D(3):315--325, 1997.

		\bibitem{bergeron2006}
		Janick Bergeron, Eduard Cerny, Alan Hunter, Andrew Nightingale,
		\textit{Verification Methodology Manual for SystemVerilog},
		Springer, New York, 2006.

		\bibitem{kroening2016}
		Daniel Kroening, Ofer Strichman,
		\textit{Decision Procedures: An Algorithmic Point of View},
		2nd edition, Springer, Berlin, 2016.

	\end{thebibliography}
	
\end{document}